\chapter{Conclusions}
\label{chap:conclusion}

This thesis set out from a structural gap in the public Telegram record: existing corpora, however large, are snapshots that record each channel once and discard almost everything that makes a channel a living object over time, its edits, its deletions, and the drift of its metadata. The resolution is a scraper that revisits public channels on a schedule, recurses into their linked discussion groups, and on each visit detects edits and deletions and snapshots channel metadata, all within Telegram's rate limits. The corpus it produces, TGDataset2, is the first general-purpose Telegram dataset to carry message edits with their before-and-after text, deletions, and longitudinal metadata, alongside a structured engagement-and-comment surface. That capability is bought with a deliberate trade of breadth for depth: $12{,}065$ channels each revisited over time.

Across the new layers, a widely used raw signal repeatedly overstates the phenomenon it is taken to measure, while longitudinal or structured capture supplies the corrected picture. The \texttt{edit\_date} timestamp over-counts human editing by roughly an order of magnitude, because $93\%$ of the edits caught in full leave the message text byte-identical, and only the before/after archive separates genuine rewrites from media and markup churn. The raw forward flag overstates editorial reposting by about a third, because a third of flagged forwards are auto-mirror copies of a channel's own posts into its discussion group; removing them moves the forward ratio from $18.3\%$ to $12.1\%$. The reply-pointer count overstates comments-on-posts by roughly $42\%$, because not every reply anchors on a broadcast post, and anchoring-aware accounting is what separates a true comment section from in-group chat.

%======================================================================
\section{Limitations}
\label{sec:limitations}

This section consolidates the known limitations of the dataset and the system that produced it. Each item is cross-referenced to the methodology section where the corresponding design decision and its rationale are discussed, so a reader can trace every limitation back to the trade-off that produced it.

\paragraph{Sampling bias from forward-only discovery.}
Channel discovery relies exclusively on forwarded-message references; username resolution is disabled to preserve the rate-limit budget for message collection (Section~\ref{sec:channel-discovery}). Channels that are only reachable through \texttt{@}-mentions or \texttt{t.me/} links, without ever being forwarded, are therefore under-sampled. More broadly, snowball sampling, the chain-referral technique that expands a seed set by following inter-channel references, is known to over-represent high-degree nodes and under-cover weakly connected regions of the network~\cite{wagner2017}. This bias is not merely theoretical here: as Section~\ref{sec:res-frontier} shows, the unreachable frontier is organised rather than random (it clusters by theme and includes the linked discussion groups of major channels), and roughly $14.5\%$ of cited \texttt{t.me} links are invite-only or web-preview references that the scraper structurally cannot resolve, a hard ceiling rather than a backlog. Any finding derived from the dataset should be interpreted in light of these biases.

\paragraph{Count-based deletion heuristic.}
The combined pass decides whether to perform a full-history deletion scan by comparing the server-reported message count against the locally stored count (Section~\ref{sec:combined-pass}). If the number of deletions between two visits is exactly offset by an equal number of new messages, the counts balance and the deletion is silently missed until a subsequent visit restores the imbalance. The full alternative, a history pull on every visit, was rejected as infeasible within Telegram's rate limits. A consequence specific to the analysis is that the deletion layer is also a detection-time record: a deleted row's timestamp marks when the scraper noticed the deletion on revisit, not when it happened, and the deletion-detection window is short (it begins in May~2026), so all deletion findings characterise the collected deleted set rather than Telegram's deletion semantics in general.

\paragraph{Edit-detection blind spot beyond the overlap window.}
On clean runs the combined pass takes Path~B (Section~\ref{sec:combined-pass}), which only re-fetches the last $N$ messages below \texttt{highest\_id} (default $N=100$) to check for edits. Edits to messages older than $N$ are therefore invisible until the same channel happens to enter Path~A for an unrelated reason: a detected deletion or a crash-triggered forced recheck (Section~\ref{sec:robustness}). These three mechanisms guarantee that no edit remains hidden indefinitely, but the latency until discovery is not bounded by $N$ alone. As a result the before/after edit archive is a selected slice, weighted toward old posts caught changing on re-scrape, and \texttt{edit\_date} remains only an upper bound on human editing (Section~\ref{sec:res-edits}). Increasing $N$ widens the window at linear API-call cost per channel per run.

\paragraph{Posted-and-deleted between visits.}
A message that is both posted and deleted in the interval between two of our visits is invisible by construction: we never observed it, the database has no row to flag, and \texttt{server\_total} does not reflect it either. This is an unavoidable property of any polling scraper and is independent of the detection strategy.

\paragraph{Race between the count ping and the new-message pull.}
The cheap path (Section~\ref{sec:combined-pass}) reads \texttt{server\_total} first and then fetches new messages. A post that lands in the few seconds between those two operations is counted in $n_{\mathrm{new}}$ but not in \texttt{server\_total}, inflating $n_{\mathrm{expected}}$ and producing a false-positive deletion signal; symmetrically, a deletion in the same window can produce a confusing comparison. The consequence is at worst one redundant Path~A run, never data loss or corruption: the full pull finds no missing IDs and the next run sees consistent counts. The races are rare enough that no mitigation is implemented. Alignment runs (Section~\ref{sec:alignment}) widen this race from the gap between two API calls to the gap between the alignment trigger and each channel's turn (minutes to hours), and change its direction: because \texttt{server\_total} remains live-global while $n_{\mathrm{new}}$ is bounded by the cutoff $t_c$, post-cutoff activity can \emph{mask} deletions rather than produce false positives. The consequence is still bounded: the regular scrape that follows every alignment run sees no cutoff, runs an exact cheap check, and catches any deletions that were masked during alignment.

\paragraph{Temporal drift across channels.}
Because channels are scraped sequentially by a single process, the per-channel metadata snapshots within a single run are captured at different wall-clock moments. The alignment snapshot protocol (Section~\ref{sec:alignment}) guarantees that every channel's \emph{message} stream is captured up to the same cutoff $t_c$, but the \emph{metadata} snapshots remain spread across the alignment cycle. Drift is therefore bounded by the alignment interval ($30$ days by default) rather than eliminated. The practical effect is visible in the analysis: the consecutive-snapshot change rates of Section~\ref{sec:res-nulls} are constrained as much by snapshot density as by genuine stability.

\paragraph{Bounded crash-window data loss.}
Messages are committed to the database per batch rather than per message (Section~\ref{sec:robustness}). A crash at the worst possible moment can therefore lose the in-flight batch (up to $100$ messages per affected channel). The lost messages are re-fetched on the next run because they sit above the stored \texttt{highest\_id}, so the loss is bounded and self-healing rather than permanent.

\paragraph{Dormant trust flag.}
The \texttt{is\_scam} flag should not be used as-is by downstream reusers (Section~\ref{sec:res-nulls}): it is effectively dead (five channels, never flipped), so any moderation-over-time analysis should rely on \texttt{is\_fake}, \texttt{is\_restricted}, or verification.

\paragraph{Companion database is separate and shallow.}
The official-Israel arm of the war showcase (Section~\ref{sec:res-war}) lives in a separate companion database, seeded from official Israeli channels and expanded by a single snowball pass, rather than in the main corpus. Cross-side comparisons are therefore reconciled in analysis on small cohorts, and they are cross-sampling-method (official-seeded voices against a subset of TGDataset's channels) as well as cross-database, so they are read as institution-fixed contrasts and per-channel shapes rather than as absolute cross-side volumes. The companion database is also shallow in time, with little edit or deletion depth, so the Israeli arm uses volume, cadence, and reactions only.

\paragraph{Scope boundaries.}
Three deliberate scope boundaries further narrow what the dataset captures. Media files attached to messages are not downloaded (Section~\ref{sec:schema}); only the media-type indicator is preserved. The scraper is designed to run as a single process on a single machine; multi-instance or distributed scraping is out of scope (Section~\ref{sec:robustness}). The unit of analysis is public channels only; private chats, groups, and user-level messages are not collected.

\paragraph{Measurement validity of derived labels.}
The analytical layers lean on off-the-shelf models applied largely cross-lingually, without task-specific ground-truth validation on this corpus: the \texttt{textdetox} XLM-R toxicity classifier~\cite{dementieva2024}, Claude Sonnet~4.6 edit labelling~\cite{claudesonnet2026}, a multilingual named-entity tagger~\cite{tedeschi2021}, BERTopic~\cite{grootendorst2022}, language detection, and cross-lingual sentence embeddings~\cite{reimers2020}, together with reaction valence and turnout read as audience-sentiment proxies. The caveat is that such results carry direction and contrast rather than calibrated absolute levels, and the firmest cross-side comparison holds the language fixed, as the war section does when it reads audience toxicity off both sides' shared Arabic sections rather than across Farsi and Hebrew.

\paragraph{Corpus composition and generalisability.}
Beyond the forward-only discovery bias above, the corpus is seeded from TGDataset's channels (Section~\ref{sec:deployment}) and is Russophone-heavy (Russian is $52\%$ of majority-language channels), so it is a particular slice of Telegram rather than a representative sample. Every rate reported in this thesis, deletion, forward, restriction, and edit, is therefore a within-corpus rate, not a platform-wide prevalence estimate, and should be read as such.

%======================================================================
\section{Future Work}

The system and corpus open more avenues than this thesis closes. The richest of them are analyses the new layers now make directly possible.

\begin{itemize}
  \item \textbf{Longitudinal moderation at scale.} The deletion-detection window only opened in May~2026, so the $2.42$\,M detected deletions are a detection-time sample rather than a record of moderation dynamics. Sustained running turns that sample into genuine deletion \emph{dynamics}, who deletes, when, and how removals propagate, and lets restriction flags and metadata be read as time series. This is the direct route from the near-null longitudinal-metadata result and the partial deletion sample of this thesis to measurable trajectories.
  \item \textbf{Edits as a research signal.} The before/after archive (currently $36{,}016$ edits caught in full) is the diffable layer no other general-purpose corpus holds. Scaled up, it supports the study of misinformation correction, stealth edits, and coordinated narrative changes from the text of the change itself rather than from a bare timestamp.
  \item \textbf{The comment layer as a social object.} With $333.2$\,M discussion messages, the comment side can carry commenter-network, returning-audience, and thread-structure studies. The anchoring-aware accounting of Section~\ref{sec:res-layers} is what makes them tractable, separating genuine comment sections from in-group chat before any measurement begins.
\end{itemize}

A smaller set of system extensions would unlock and sharpen those analyses. Lower-latency edit and deletion detection (an adaptive overlap window, selective history pulls, and a limited real-time API path in the spirit of Kireev et al.~\cite{propaganda2025}) would catch messages posted and deleted between visits and recover true deletion timestamps rather than detection times; denser, event-triggered metadata snapshots would resolve the temporal-drift limitation that holds the current metadata result near null; re-enabled username and link discovery under a separate rate budget would shrink the organised frontier; optional media or perceptual-hash capture would let the sockpuppet null result be revisited on image content rather than profile pictures alone.

%======================================================================

Making content \emph{dynamics} observable, not just content, is the shift the title names, and it is what changes the questions one can put to Telegram: from what a channel said once to how its public record was edited, deleted, and otherwise managed over time. That the system does this openly and reproducibly from the documented defaults of Table~\ref{tab:config} is the point. It lowers the barrier for others to ask those questions of the platform too.
